% \section*{About the coursework}
\chapter{Introduction}

\section{About the Assignment}

\subsection{Introduction}

Product development is a multifaceted endeavor that relies on contributions from various functions within a company. To reflect the actual product development activity, the aim of this group project is twofold.

\begin{enumerate}
    \item To reflect the interdisciplinary nature of product development.
    \item To empower individuals through experiential learning and foster a culture of learning by doing.
\end{enumerate}

Regarding the first aim, the weekly assignment (i.e., tutorial session) will provide you with valuable experience in product development as well as exposure to different aspects of teamwork. Regarding the product development exposure, this mean you will experience activities involving identifying opportunities, planning products, and understanding customer needs as a team. Whereas, for the exposure to different aspect of activities, this mean, you will involves in conducting meetings, making decisions, managing time and people, writing reports, delivering presentations, improving interpersonal skills, and resolving conflicts. You will also discover that working on product development as a group can be very enjoyable.

As for the second objective, I strongly believe that learning by doing facilitates better retention of knowledge and enhances students' ability to absorb and comprehend concepts.


\subsection{Objective}

The objective of this group work assignment is to design a product of your choice (consensus of your team) and provide a description of the design and development process.



\subsection{Product Selection}

As mentioned earlier on, the objective of this group coursework is  to help you understand each session's topic better, this APID course asks you to work in groups and together, you'll think of an idea, design it, and present a prototype or pitch the idea to public tentatively at week 14 (I hope to made it this semester). The ideas for projects will come from you, the students.

This semester mark the third time this course is being offered. Previously, I always encourage students to select a product that they are interested in, and I have been very impressed with the creativity and diversity of the products that students have come up with. However, I have noticed that some students have struggled to come up with a product idea, and this has caused some stress and anxiety. Therefore, this semester, pick a theme for your project. The theme is "Hydroponics Agriculture". This means that your project should be related to hydroponics agriculture in some way. You can choose to design a product that is used in hydroponics agriculture, or you can design a product that is used to support hydroponics agriculture. The choice is yours, but it must be related to the theme.


\subsection{Product development}
Product development is a series of activities that starts with identifying a market opportunity and ends with the production, sale, and delivery of the product. Your product design should address all the required items. If some items are not well addressed, you will lose points.



\subsection{Group}

At the start of the module, you will be divided into groups of around four or five. The group division is done randomly. This is done to ensure diversity and to simulate the real world, where you will often work with people you don't know. Unfortunately ,you are not allowed to choose your own group or to switch groups with someone else. But, do remember that, working with people from different backgrounds is a skill that will be important in your future professional life.

Once the randomization is completed, the group allocation can be found on the module's ITEL page.

\textbf{Important!} If you are not listed as a student in the module, or if you have not been assigned to a group, you must contact the module convenor as soon as possible. This is especially important if you are an exchange student, as you may not be assigned to a group right away. If you cannot contact certain group members, such as if they have switched modules, you must also let the convenor know.


Each group should appoint a group leader. To be fair, the group leader should be rotated every one or two week. The group leader will play a management role in addition to their normal work as a group member. This includes ensuring that everyone is involved, that work is divided fairly, that tasks are completed on time, and that nothing is forgotten in the project. The group leader does not have to make all the decisions, but they do have a responsibility to ensure that the group works effectively as a whole. They should also ensure that everyone is pulling in the same direction and that the project progresses according to plan. Ultimately, the group leader is responsible for ensuring that the project goals are achieved and that the prescribed deadlines are met.
Once you are assigned to a project team, you should remain to stay in the team for the entire term.


\subsection{Meetings}
Each group should arrange meetings at their own discretion. The meetings will normally be held in a mutually convenient location in the faculty, or elsewhere on campus, and have no prescribed time limit. What goes on at these meetings is completely up to the group, but they are usually more informal in nature. There is no requirement for a chairperson or minutes, but some groups may find it useful to carry on with these practices here.


\subsection{Pre-Class and Collaborative Task}
There are certain activities that need to be completed before each session, which will be discussed and shared with other students during each session. Additionally, some activities are designed for group work and should be carried out outside the classroom. However, specific details regarding the location and timing of these activities will be communicated or emphasized in the text.
During the pre-class and in-class discussions, the main focus is on brainstorming and expanding ideas, so accuracy is not the top priority. You only need to present your ideas in bullet point format. However, it is expected that in the final report, all details will be clarified and presented in a suitable reporting format.



\subsection{Project Materials and Expenses}
Unfortunately, we cannot provide money for your project. If your project needs extra money, your team will need to pay for these costs.


\subsection{Intellectual Property Rights}
In this course, any inventions made by your team will usually belong to you. If you decide to apply for a patent to protect your invention, you can either do it by yourself and pay all related costs, or you can choose to patent it with the help of Universiti Malaysia Sabah. In this case, the university will take care of all expenses and other necessary steps. However, it's important to discuss and agree with your team members early on about how to divide any earnings from your invention. Also, any ideas shared on a digital whiteboard or similar tool should be considered as recorded evidence of how your invention developed over time.







\subsection{A Few More Hints}
\begin{itemize}
    \item Keep your highly potential/ proprietary ideas for another context. In our class, we'll talk openly about the projects, and we want to avoid being limited by secret information.
    \item Projects often do better when at least one member in the team is really interested in what the project is about.
    \item Just because you've come across a product that doesn't work well doesn't mean there isn't a better version out there. Make sure to look around carefully to see what else is available.
\end{itemize}


